\documentclass[12pt]{letter}
\usepackage[margin=1in]{geometry}
\usepackage{hyperref}

\signature{Quang-Vinh Dang}
\address{Quang-Vinh Dang \\ British University Vietnam \\ Hung Yen, Vietnam \\ vinh.dq4@buv.edu.vn}

\begin{document}

\begin{letter}{Editor-in-Chief \\ Electronic Commerce Research and Applications \\ Elsevier}

\opening{Dear Editor,}

I am pleased to submit my manuscript, ``When AI Shops with AI: Buyer-Seller Agent Negotiation and Market Outcomes in Agentic Commerce,'' for consideration for publication in \emph{Electronic Commerce Research and Applications}.

As AI shopping agents increasingly search, compare, and negotiate on behalf of buyers, and sellers respond with their own AI pricing and negotiation agents, the institutional rules of e-commerce markets that were once set implicitly by human cognitive and time constraints -- how much of the market gets searched, whether price is posted or bargained, what the bargaining space is allowed to include, and how concentrated the agents themselves are on shared infrastructure -- become explicit design choices with measurable welfare and competitive consequences. This paper studies those consequences using a controlled, matched-sample agent-based-computational-economics (ACE) simulation of 105,600 transactions across 80 replicate markets, crossing five institutional regimes with seller information, bargaining scope, and buyer- and seller-agent-provider concentration, and corroborates its central finding with a smaller, exploratory pilot in which real large language models negotiate directly.

The paper's central contributions are: (1) a clean decomposition showing that AI-mediated search improves welfare through better matching, that price-only bargaining is close to purely redistributive, and that only expanding the bargaining space to include bundles and terms creates genuine additional surplus; (2) evidence that richer seller information about buyers, a natural byproduct of AI personalization, systematically compresses the share of that surplus captured by buyers; (3) a demonstration that AI-mediated search sharply concentrates trade among sellers, an effect that competitive multi-seller solicitation meaningfully but only partially reverses; and (4) a careful empirical and conceptual distinction between price correlation among sellers sharing third-party agent infrastructure and genuine algorithmic collusion, with direct implications for how competition regulators might approach agentic commerce. Throughout, the paper is explicit about the difference between what a rule-based simulation can and cannot claim to show, and reports a set of robustness checks -- including a sensitivity analysis on the paper's key structural parameter and a buyer-fixed-effects specification exploiting its matched design -- added specifically to test the paper's own assumptions rather than merely assert their robustness.

This manuscript is original, has not been published previously, and is not under consideration at any other journal. I have no conflicts of interest to disclose. I believe the paper's combination of institutional-design questions, competition-policy relevance, and methodological transparency about the scope of a rule-based simulation makes it a good fit for \emph{Electronic Commerce Research and Applications}'s readership.

Thank you for your consideration. I look forward to your response.

\closing{Sincerely,}

\end{letter}
\end{document}
